% Lösungen Einheit 2 von 3 – Rate und Bestand als Paar (zusammenbau v0.1)
\einheitenkopf{Einheit 2 von 3 · Rate und Bestand als Paar}

\erg{11}{a) $t = 5$ – dort wechselt r von plus nach minus \quad b) Nullstellen von r: $t = 0$ und $t = 8$; dazwischen ist $r(t) > 0$ (etwa $r(4) = 8$), also nimmt der Inhalt durchgehend zu \quad c) $t = 6$: für $2 < t < 6$ ist r positiv, danach negativ – dort wechselt r von plus nach minus. Für $0 < t < 2$ ist r negativ, aber $∫_0^6 r(t)\,\mathrm{d}t = 18 > 0$, also ist die Schlange bei $t = 6$ länger als zu Beginn \quad d) $L'(t) = 6t - 3t^2 = r(t)$ und $L(0) = 0$ (niemand zu Beginn) – beides nötig; $L(3) = 27 - 27 = 0$: aufgelöst \quad e) $W(t)$ ist die Höhenzunahme seit der Pflanzung (das Integral der Rate von 0 bis t), $W(t) + 1{,}5$ die Höhe nach t Jahren; Grenzwert $= 18$: auf lange Sicht nähert sich der Baum einer Höhe von 18\,m}
\erg{12}{a) $r(t) = 0$ für $t = 0$, $t = 5$ und $t = 10$; Testwert $r(6) = -12 < 0$: das Volumen nimmt für $5 < t < 10$ ab}
\erg{13}{a) Emma nimmt den Hochpunkt der Rate, dort wächst der Bestand nur am schnellsten. Größter Bestand, wo r von plus nach minus wechselt: $r(t) = t \cdot (8 - t) = 0$ bei $t = 8$ \quad b) Vorher ist die Rate positiv und der Bestand wächst, danach ist sie negativ und er fällt; dazwischen liegt der höchste Wert. \quad c) B steigt bis t = 4 (dort ist r null, Hochpunkt von B) und fällt danach; am steilsten steigt B bei t = 0; $B(4) = 5$, $B(6) = 4$ \quad d) r ist positiv für $0 < t < 7$ und wechselt bei $t = 7$ von plus nach minus: die meisten Gäste um 16 Uhr}

13c)

\begin{ksys}[xmin=0,xmax=7,ymin=0,ymax=6,klein] \funktionab{1+2*\x-0.25*\x^2}{B}{0}{6} \end{ksys}

